\zlabel{triptych:concise:themes:start}
\section{The Propers: Themes and Movement}
Readiness to meet Christ requires mercy toward the neighbor and charity within the act of mercy. Chrysostom asks how people distinguished by a demanding virtue can neglect someone in need. Augustine asks how even those who help can lack love. Gregory explains both the hospitality through which the Emmaus disciples recognize Christ and the desire for praise that can corrupt generous acts. The two interpretations examine different failures in lives that may appear faithful.

\subsection{Asking, seeking and receiving help}
The entrance antiphon adapts the plea of Psalm 87:3 Vulgate, 88:3 Hebrew. The whole psalm leaves its speaker in distress. Augustine hears Christ's human suffering in that prayer and understands divine hearing as merciful acceptance. Sorrow can therefore accompany a truthful approach to God. The Collect then asks for freedom directed toward service. Alongside the plea to be heard, that petition can include becoming available to another person's need; alongside Augustine's examination of purpose, it can include release from the desire to be admired.\footnote{Augustine, \textit{Expositions on the Psalms}, 87.1--4.}

In the first reading, Wisdom is loved and sought, yet she also makes herself known and goes out to meet the seeker. The final appointed verse includes her own search. Bonaventure explains that Wisdom's light makes her visible and that the intellect and will must open to her. Seeking involves action as well as thought. Read beside the Gospel's warning, this approach from Wisdom gives confidence to begin preparing now: the response asked of the seeker begins with a gift.\footnote{Bonaventure, \textit{Commentarius in librum Sapientiae}, 6:13--17 Vulgate, Quaracchi VI, pp. 147--149.}

The Responsorial Psalm describes thirst with soul and flesh, praise with lips and lifted hands, and remembrance on the bed. Augustine extends the body's desire toward its resurrection. He also connects the hands raised in prayer with hands extended to the needy. Bodily rest is compatible with the soul's vigilance. These details allow both practical care and inward attention to belong to the same response to God.\footnote{Augustine, \textit{Expositions on the Psalms}, 62.1--12, especially §10.}

\subsection{Hope for the dead and preparation for the arrival}
The second reading addresses mourners. Both forms begin with Jesus' death and resurrection; the long form continues through the rising of the dead and the gathering with the living to meet the Lord. It concludes by directing believers to console one another. Aquinas distinguishes grief over separation from hopeless sorrow and explains Christ's rising as the cause and pattern of ours. Comforting someone can therefore acknowledge the loss while giving the reason for hope. The trumpet and gathering belong to the long form and to the short form's surrounding context.\footnote{Aquinas, \textit{Super I ad Thessalonicenses}, 4, lectio 2.}

The acclamation calls for watchfulness because the hour is unknown. In Matthew's surrounding discourse, Chrysostom explains readiness through the servant's continuing responsibility for others. In the Gospel, all ten virgins become drowsy and sleep. Augustine reasons that their common sleep cannot mean charity growing cold: that would make the wise fail too. He interprets it as death. The distinction disclosed at the arrival concerns what each has prepared, rather than whether the body has rested.\footnote{Chrysostom, \textit{Homilies on Matthew}, 77.1--3; Augustine, Sermon 93, §§6--9.}
\newpage
\subsection{Mercy performed and the purpose of performing it}
Chrysostom identifies the lamps with virginity and the oil with humanity, almsgiving and assistance. An arduous discipline can coexist with attachment to money and neglect of a neighbor. He therefore asks why someone who has resisted a powerful appetite should allow a more ordinary attachment to spoil that labor. The difficulty becomes acute when the wise refuse oil: if oil means generosity, their refusal needs explanation. Chrysostom answers that a life of mercy cannot be transferred at the final encounter. Present life is the time to give.\footnote{Chrysostom, \textit{Homilies on Matthew}, 78.1--2.}

Augustine identifies the lamps with good works, including feeding the hungry and welcoming a stranger. Those actions can occur without inward charity. Gregory likewise describes almsgiving done for human favor and interprets the vessels as hearts. The foolish seek other people's testimony at an encounter where God knows the purpose of the acts. Their reliance on praise cannot provide what they lack within. This examination asks the giver to purify the purpose of service while continuing to serve.\footnote{Augustine, Sermon 93, §§2--5, 9--14; Gregory, \textit{Homiliae in Evangelia}, 12.1--3.}

The difference includes the sellers. Chrysostom identifies them with the poor, whom the foolish should have helped during life. Augustine and Gregory identify them with flatterers, whose approval cannot secure admission. Their moral demands are compatible; these explanations of the same figures differ. Both concern a preparation that must become the person's own. Neither permits withholding assistance from someone who can still be helped, and neither makes God's gift something purchased by a sum of money or secured by a favorable reputation.

\subsection{Receiving Christ and persevering after reception}
The Prayer over the Offerings concerns celebrating the Son's Passion and receiving it inwardly. In the first interpretation, receiving Christ's self-giving calls for help that accepts a cost for another person. In the second, participation includes examination of the purpose with which one acts. The relation between outward celebration and inward response allows both questions to be asked without reducing the sacrament to the participant's feelings.

The two Communion antiphons offer distinct scenes. The psalm alternative names the shepherd's provision. Augustine identifies Christ caring for the Church and leading to baptismal refreshment. Bellarmine explains that divine care can coexist with hardship. Gratitude for receiving care can encourage care for others, even where material need persists. In the Luke alternative, the disciples report their recognition of Jesus in breaking bread. Gregory explains how receiving the unrecognized traveler precedes understanding. They act on the word before grasping who their guest is.\footnote{Augustine, \textit{Expositions on the Psalms}, 22.1--6; Bellarmine, \textit{Commentary on the Psalms}, O'Sullivan abridgment, Psalm 22; Gregory, \textit{Homiliae in Evangelia}, 23.1--2.}

After Communion, thanksgiving is followed by a petition for the Spirit to preserve integrity. Mercy must continue when helping is inconvenient; charity must endure when good work receives no praise. The petition entrusts that continuing response to God's help. The promised encounter with Christ includes both his judgment and the lasting life with him of which Paul speaks.
\zlabel{triptych:concise:themes:end}
\clearpage
